\chapter{BCS Theory of Superconductivity}

Superconductivity was experimentally observed well before the formal discovery of quantum mechanics, presenting a profound physical mystery that eventually necessitated the development of a comprehensive quantum many-body description. The definitive microscopic explanation of this phenomenon, known as the BCS theory, was developed in 1957 by John Bardeen, Leon Cooper, and John Robert Schrieffer. At its core, the theory explains the superconducting state through the formation of bound states of electrons, known as Cooper pairs.

To formally develop this, we will consider a general system composed of spin-1/2 fermions experiencing a mutual \emph{attractive} interaction between particles of opposite spin. Our primary theoretical interest is to determine the low-temperature phase of this many-body system. The central idea is that this mutual attraction can induce the formation of Cooper pairs—pairs of fermions with opposite spin. Because these composite pairs possess an integer total spin, they behave effectively as bosons and can undergo macroscopic condensation at sufficiently low temperatures.

The original motivation for studying such an interacting fermionic system was the microscopic explanation of metallic superconductivity. Superconductivity is a physical phenomenon occurring in certain metals and compounds whereby an electric current flows with absolutely zero electrical resistance below a specific critical temperature \( T_c \). This phenomenon is deeply analogous to superfluidity; whereas a superfluid involves the macroscopic flow of a fluid with zero viscosity (the mechanical analog of electrical resistance) below a critical temperature, a superconductor involves the frictionless flow of charged carriers, namely electrons.

A crucial thermodynamic clue pointing toward the existence of these bound pairs of fermions in a superconductor is provided by the electronic specific heat, \( c_v \). For an ideal Fermi gas at low temperatures \( T \)—which serves as a highly effective model for the conduction electron gas in a normal solid—the specific heat scales linearly with \( T \). However, when a material transitions into a superconductor, the specific heat is observed to vanish exponentially with decreasing temperature. This exponential suppression is the hallmark thermodynamic behavior expected when there is a finite energy gap in the excitation spectrum separating the macroscopic ground state from the first available excited states. The magnitude of this energy gap can be extracted directly from specific heat measurements, revealing it to be extremely small—typically several orders of magnitude smaller—compared to the Fermi energy \( E_F \) of the electrons. It is natural to guess that this energy gap is directly associated with the finite binding energy of the fermion pairs. (The microscopic mechanism by which normally repulsive electrons can experience an effective mutual attraction in a metal lattice will be discussed in detail in the subsequent section).

Beyond its foundational success in condensed matter systems, the BCS theoretical framework we are going to develop is remarkably universal. The collection of phenomena that can be governed by these same equations spans an astonishing range of more than ten orders of magnitude in temperature. At the lower extreme, in the nanokelvin regime, the theory applies directly to ultracold Fermi gases trapped in optical lattices. These are dilute, neutral systems of fermionic atoms where an effective mutual attraction—arising fundamentally from van der Waals forces—can be accurately modeled as a simple contact interaction due to the extreme diluteness of the gas. These highly tunable systems have been studied very actively over recent years because experimentalists can artificially tune the strength of the interparticle interaction at will. The theory is equally relevant to other macroscopic neutral systems, such as liquid \( ^3\text{He} \), where fermionic neutral atoms also experience a mutual attraction. At the absolute opposite extreme of the temperature and density scale, the very same BCS formalism can be extended to describe the pairing of quarks inside the ultra-dense, extremely high-temperature cores of neutron stars (a phenomenon known as color superconductivity).

\section{Basic Model for Effective Attraction}

Let us consider a simplified model of a solid consisting of a cubic lattice of positively charged ions with a lattice spacing \( a \), permeated by a gas of itinerant electrons. We assume the system is isolated, without any external electric or magnetic fields. The electrons are subject to a periodic potential generated by the ionic lattice; consequently, the potential minima are located in the immediate vicinity of the positive ions, marking the regions where the electrons are most likely to be found.

The electrons interact not only with this periodic ionic potential but also with one another via the standard Coulomb repulsion, \( V_{e-e} \sim e^2/r \). However, the presence of the mobile electron gas embedded in the positively charged background activates the many-body phenomenon of screening. In a metal, the spatial redistribution of conduction electrons around any single negative charge effectively screens that negative charge at a distance from it. This screening effect drastically modifies the long-range Coulomb repulsion into a short-range effective interaction, which decays exponentially as \( \sim e^{-r/\xi} \), where the screening length \( \xi \) is typically on the order of a few lattice spacings \( a \). Although this screened interaction is significantly weaker at macroscopic distances, it remains fundamentally repulsive. To explain the formation of Cooper pairs, we must identify a physical mechanism capable of overcoming this residual repulsion to produce a net effective attraction between the electrons.

\subsection{The Electron-Phonon Interaction and Fröhlich Mechanism}

This requisite attraction arises from the interaction between the electrons and the collective vibrations of the ionic lattice, known as phonons. Because the typical Debye temperature characterizing these macroscopic phonon modes is on the order of \( \omega_D \sim 100\,\text{K} \), while superconducting phenomena typically occur at cryogenic temperatures around \( 1\,\text{K} \), the thermal phonons are effectively frozen out. Consequently, we can safely treat the quantized phonon field as a perturbation to the interacting electron system.

As famously demonstrated by H. Fröhlich, in certain metals, the combination of the screening of the Coulomb repulsion and the virtual exchange of phonons between electrons may lead to an overall effective attraction between them. Specifically, the dynamic interaction between electrons and phonons produces a retarded attractive force. When an electron moves through the crystal, it locally modifies the ionic lattice by pulling the positive ions slightly toward its trajectory, creating a localized cloud of positive charges, and then rapidly moves away. Subsequently, a second electron may be electrostatically attracted by this localized accumulation of positive charge.

A key physical ingredient of this retarded mechanism is the stark difference in dynamical time scales: the heavy distortions of the ionic lattice fade away rather slowly in comparison with the extremely fast time scale over which one electron is replaced by a second one at a certain position due to the rapid motion of the electrons. Ultimately, these two combined effects—electronic screening and virtual phonon exchange—may lead in certain materials to an overall effective attraction between the electrons, provided those electrons belong to a tiny energy shell situated immediately around the Fermi energy \( E_F \).
